\section{Composed Duality Action and Observable Constraints}
\label{sec:composed_duality}

\paragraph{Composed representative action.}
Define
\begin{equation}
  \mathcal{D}_{\mathrm{comp}} := \mathcal{G} \circ \mathcal{D}_{\mathrm{sw}}.
\end{equation}

\begin{proposition}[Composed-action structure]
\label{prop:composed_action}
The maps \(\mathcal{D}_{\mathrm{sw}}\) and \(\mathcal{G}\) are both involutions on representative
labels. Their action on the radial sector is always commuting:
\begin{equation}
  (\mathcal{G}\circ\mathcal{D}_{\mathrm{sw}})\chi
  =(\mathcal{D}_{\mathrm{sw}}\circ\mathcal{G})\chi
  =\chi^{-1}.
\end{equation}
On oriented sectors, commutation holds whenever the dual map preserves orientation parity
\(u^\vee(-u)=-u^\vee(u)\). In that case, the generated discrete group is \(Z_2\times Z_2\).
\end{proposition}
\begin{proof}
Involution of \(\mathcal{D}_{\mathrm{sw}}\) follows from Theorem~\ref{thm:strong_weak_involution},
and involution of \(\mathcal{G}\) is immediate from sign flip. The radial statement follows because
\(\mathcal{G}\) acts trivially on \(\chi\). Parity preservation is the sufficient condition for commuting
action on \(u\).
\end{proof}

\paragraph{Constraint on representative fits.}
Suppose a saturated-regime radius is parameterised as \(r=r(\beff)\). Orientation pairing then
imposes
\begin{equation}
  r' = r(2\beta - \beff).
\end{equation}
This gives a non-trivial consistency relation for any empirical \(\beff\)-collapse.

\begin{corollary}[Ratio constraint]
\label{cor:ratio_constraint}
\begin{equation}
  R(\beta,\theta) := \frac{r(2\beta-\beff)}{r(\beff)}.
\end{equation}
At the symmetry point \(\Theta=0\) (equivalently \(\beff=0\)), the paired representatives are
\((0,2\beta)\), and \(R\) provides a sharp local diagnostic of representative consistency.
\end{corollary}

\paragraph{Temperature-driven branch exchange.}
The dual map also yields a direct statement about high-\(T\)/low-\(T\) control transfer when
\(\chi_0\) has monotone endpoint scaling.

\begin{corollary}[Classical-quantum exchange of control]
\label{cor:classical_quantum_exchange}
Assume one-channel scaling \(\chi(\beta,g,\theta)=g^2\chi_0(\beta,\theta)\) with \(g>0\), and
suppose
\begin{equation}
  \chi_0(\beta,\theta)\to 0 \quad (\beta\to 0),
  \qquad
  \chi_0(\beta,\theta)\to \infty \quad (\beta\to\infty).
\end{equation}
Then, at fixed \(g\),
\begin{align}
  \beta\to 0 &: \chi(\beta,g,\theta)\to 0, \qquad g^\vee(\beta,\theta;g)\to\infty, \\
  \beta\to\infty &: \chi(\beta,g,\theta)\to\infty, \qquad g^\vee(\beta,\theta;g)\to 0.
\end{align}
Hence low-temperature strong-\(\chi\) asymptotics at \(g\) are transferred to weak-\(\chi\)
asymptotics at \(g^\vee\), while high-temperature weak-\(\chi\) asymptotics map to a
strong-coupling dual representative.
\end{corollary}
\begin{proof}
From Theorem~\ref{thm:strong_weak_involution},
\(g^\vee(\beta,\theta;g)=1/(\chi_0(\beta,\theta)g)\) and
\(\chi(\beta,g^\vee,\theta)=\chi(\beta,g,\theta)^{-1}\). The endpoint limits follow by
substitution of the assumed \(\chi_0\) scaling.
\end{proof}
